\chapter{Capítulo Primero - Introducción }
La presente memoria recoge el trabajo realizado en el marco del Proyecto de Fin de Grado del Grado en Ingeniería Informática. Este proyecto integra de forma aplicada los conocimientos y competencias adquiridas a lo largo de la titulación, especialmente en las áreas de ingeniería del software, arquitecturas en la nube, gestión y procesamiento de datos, sistemas inteligentes y análisis predictivo.

El proyecto desarrollado consiste en la evolución del prototipo académico \textit{TankGo} \cite{Alvis2025TankGo} hacia una plataforma software completa orientada a la consulta, comparación y análisis de precios de combustibles y puntos de recarga eléctrica. La solución propuesta se concibe como un sistema integrado desplegado en entorno cloud que combina distintos componentes tecnológicos: ingestión y persistencia de datos, procesamiento distribuido, modelos de predicción basados en aprendizaje automático, algoritmos de recomendación geoespacial y un sistema de asistencia conversacional para la interacción con el usuario.

Este proyecto se enmarca en el área de Sistemas de Información e Ingeniería del Software, al abordar el diseño, desarrollo, despliegue y evaluación de una solución tecnológica orientada a resolver un problema real mediante una arquitectura escalable y modular.

\section{Motivación}
El sector de las gasolineras y los puntos de recarga eléctrica presenta una alta variabilidad en precios y disponibilidad, lo que genera incertidumbre en los usuarios y dificulta la toma de decisiones. Aunque existen herramientas que permiten consultar precios actuales, pocas soluciones integran funcionalidades avanzadas como predicción a corto plazo, recomendación basada en rutas o análisis geoespacial optimizado dentro de una misma plataforma.

La motivación principal de este proyecto es explorar la viabilidad técnica de integrar en un único sistema distintos enfoques tecnológicos tales como analítica de datos, predicción, recomendación y asistencia inteligente con el fin de aportar un valor respecto a soluciones convencionales.

El proyecto es una evolución de un desarrollo académico previo, ampliando su alcance, mejorando la arquitectura y profundizando en la fundamentación técnica y metodológica exigida en un Proyecto Fin de Grado.

\section{Estructura de la memoria}
La memoria se ha estructurado en capítulos que recogen de manera ordenada los distintos aspectos del proyecto. Tras esta introducción, se presentan los antecedentes y la justificación técnica del sistema, incluyendo el estado del arte y el análisis de soluciones existentes.

Posteriormente, se definen los objetivos y el alcance del proyecto, seguidos de la planificación temporal y de recursos, el presupuesto estimado y la metodología adoptada.

El núcleo del documento lo constituye el capítulo de desarrollo, donde se detallan la especificación de requisitos, el diseño del sistema, las decisiones tecnológicas adoptadas, la implementación y el plan de pruebas. Finalmente, se incluye una valoración ética del proyecto, las conclusiones obtenidas y posibles líneas futuras de mejora, junto con la bibliografía y los anexos correspondientes.